
\chapter{Background and Purpose}

\section{Background}
Lunar rovers are mobile vehicles designed to be used on the Moon's surface. They perform tasks such as terrain analysis, surface mapping, image collection, and measurements.
Rovers provide surface data that can be analysed in detail after collection.

The first successful lunar rover was the Soviet Lunokhod-1, which landed on November 17, 1970, and operated for 11 months,
transmitting thousands of images and videos back to Earth \cite{FirstRover}. Since then, rovers have become an important tool for planetary exploration.
NASA's Mars rovers (Spirit, Opportunity, Curiosity, Perseverance) have shown complex mobility and autonomous navigation capabilities \cite{nasa_mars_rovers}.


\section{Problem Statement}
Real planetary rover missions are expensive and highly specialized, which makes full-scale rover development unsuitable for most student projects \cite{perseverance_cost}.
Implementing a full-scale rover is difficult within the scope of a student project. Therefore, this project investigates whether a low-cost prototype can demonstrate selected core rover functions.

The prototype focuses on the following core functions:
\begin{enumerate}
    \item Remote operation on simulated lunar terrain.
    \item Pseudo-3D mapping using a rotating 2D LiDAR.
    \item Closed-loop stepper motor control for drivetrain and LiDAR rotation.
    \item Low-light operation without relying on camera-based navigation.
    \item ROS 2-based software integration.
\end{enumerate}

\newpage
\section{Purpose and Objectives}
The project was originally inspired by the LOOK competition held in Kongsberg, Norway. As the project developed, our group focused on a self-defined rover prototype for remote operation and pseudo-3D mapping in a simulated lunar environment.
The competition, organized in collaboration with the Norwegian Space Cluster, challenges teams to develop solutions for lunar rover problems \cite{look}.


Our specific objectives were inspired by the LOOK competition context, and are as follows:
\begin{itemize}
    \item \textbf{Remote operation}: Control the rover through a program running locally.
    \item \textbf{Mobility}: Design a rocker-bogie-inspired suspension and differential steering system capable of driving on uneven terrain with small rocks and loose regolith.
    \item \textbf{Mapping}: Generate a pseudo-3D point cloud using a 2D LiDAR rotated by a closed-loop stepper motor, enabling low-light mapping.
    \item \textbf{Power management}: Operate on lithium-ion battery power with voltage and current monitoring.
    \item \textbf{Sensor integration}: Use data from wheel encoders and an IMU for position estimation.
    \item \textbf{Size limit}: Needs to be within 0.5 m x 0.2 m x 0.2 m fully assembled
    \item \textbf{Weight}: Maximum target: 15 kg. This limited material use, motor count and chassis size.
\end{itemize}

\section{Scope}
\textbf{Within scope:}
\begin{itemize}
    \item Remote control via local software.
    \item Differential steering with closed-loop stepper motors.
    \item 2D LiDAR with rotating mount for pseudo-3D mapping, Point cloud visualization.
    \item IMU-based heading for odometry.
    \item ROS 2 software.
\end{itemize}
